% ------------------------------------------------------------------------------
% Chapter 4
% ------------------------------------------------------------------------------
\chapter{Result} 
\label{cha:result}

To covert the classical testbench into a UVM one we modified the general structure of some modules, created new ones, introduced the coverage and added some features to make testing smooth.

\section{Changes}

Here we will introduce all the logic changes from the classical testbench to the UVM one.
The changes will be explored in the following subsections.

\subsection*{Transaction (sequence item)}

The transaction used in the initial testbench wraps both masters and some unuseful signals such as \texttt{use\_m*} and \texttt{is\_end}.
These informations will be handled by the UVM structure and the driver idle state.
The final product is in charge of only one master and leaves the burden of the two masters conflicts to the sequences.
There is a second (outputs) and third (reset) transactions that handle the communication between monitors and scoreboard.

\subsection*{Virtual Interface \& Clocking Block}

Signals will be abstracted into an \texttt{interface} and synchronized with a clocking block (\texttt{cb}).
Following the transaction division of signals into masters, the \texttt{interface} abstract the control of a single master, this will lead to a clearer communication between UVM components.
This element is instantiated in the top module, and distributed through the \texttt{uvm\_config\_db}.

\lstinputlisting[
    style=svlog, 
    linerange={12 - 30}, 
    firstnumber=12
]{controller/tb_uvm_simple_mem_ctrl.sv}

\subsection*{Generation (sequence / sequencer)}

The generation of transaction was handled with a repetitive function, hard to scale and to verify.
Following the UVM architecture this will be handled by sequences that are easy to use, and to change if necessary.
Furthermore, the clocking collisions are handled by the UVM sequencer of each agent, so no manual debugging will be necessary.
There are two main tests:

\begin{itemize}
    \item \textbf{Deterministic} tests are known before the execution and are used to test a specific feature:
    \begin{itemize}
        \item \textbf{Single master write + read}
        \item \textbf{Contention tests} where both masters request the memory at the same time
    \end{itemize}
    \item \textbf{Random} tests are used to test the DUT with some general behaviors
\end{itemize}

\subsection*{Driver}

The original module took care of both master at the same time, added clock delays and applied the transactions to the DUT signals.
The new driver is compartmentalized to only one master and does not wait useless cycles that would make back to back instruction impossible.
It also handles reset sequences to not lose transactions.

The key element of this structure is the \texttt{apply} task.
This function applies the delay of each transaction, wait for the ack signal, and leaves the enable signal low in case there is no next instruction to be executed.
This last feature is a key feature to make the two masters driver work simultaneously without creation of false commands and spin locks.

\lstinputlisting[
    style=svlog, 
    linerange={105 - 116}, 
    firstnumber=105
]{controller/tb_uvm_simple_mem_ctrl.sv}

\subsection*{Monitor}

The key difference in the monitors is the use of a specialized transaction that made communication straight forward.
Having multiple monitors (one for each master + one for the reset) added complexity to the scoreboard model but UVM structures made this a minor problem.

\subsection*{Scoreboard}

In terms of implementation there is a big difference, but the main logic stayed the same.
The only real difference is the ability to manage the reset signal and the duty to stop the simulation when the last transaction is processed.

\section{Added}

The UVM architecture needed some structural blocks and to get a complete testbench some new features were added.
The added component will be explored in the following subsections.

\subsection*{Test, Environment, Agent}

To complete the structure of UVM these structures were added, they are the node of the tree of modules that connect all the other components.
This creates a modular code that enables easy changes if necessary, and makes adding new features fast and uncomplicated.

\subsection*{Coverage + Checking}

To complete the testbench all the necessary coverage points have been added: from the hardware coverage to the scoreboard coverage.
A checker was also added to check the hardware wires behavior.